%% kapa: exam=1 round=36 year=2025 period=1 form=A subject=real_estate subject_ko=부동산학원론 seq=25 number=105
%% source: Q-Net 감정평가사 기출문제 (cst00309) artlSeq=5231525 file=raw/1008/5231525/2025년 제36회 감정평가사 1차시험 1교시 문제지 원본.pdf
%% conversion: kordoc 4.24.51 -> tools/gen1.py (auto); review=proofread-sonnet; flags=none
\begin{question}{105}
국토의 계획 및 이용에 관한 법령상 기반시설의 유형으로 옳지 않은 것은?
\begin{choices}
  \item 공공ㆍ문화체육시설: 광장ㆍ공원ㆍ녹지ㆍ유원지ㆍ공공공지
  \item 유통ㆍ공급시설: 유통업무설비, 수도ㆍ전기ㆍ가스ㆍ열공급설비, 방송ㆍ통신시설, 공동구ㆍ시장, 유류저장 및 송유설비
  \item 보건위생시설: 장사시설ㆍ도축장ㆍ종합의료시설
  \item 방재시설: 하천ㆍ유수지ㆍ저수지ㆍ방화설비ㆍ방풍설비ㆍ방수설비ㆍ사방설비ㆍ방조설비
  \item 교통시설: 도로ㆍ철도ㆍ항만ㆍ공항ㆍ주차장ㆍ자동차정류장ㆍ궤도ㆍ차량 검사 및 면허 시설
\end{choices}
\end{question}
\answer{1}
